% expose_main.tex — Exposé Masterarbeit Joschka Peeters
% Thema: Living Knowledge Graph für Coding Agents

\documentclass[12pt,a4paper]{scrartcl}
\usepackage[ngerman]{babel}
\usepackage[utf8]{inputenc}
\usepackage[T1]{fontenc}
\usepackage{lmodern}
\usepackage[style=authoryear-comp,backend=biber,maxcitenames=2,maxbibnames=6]{biblatex}
\usepackage{csquotes}
\usepackage{geometry}
\usepackage{setspace}
\usepackage{microtype}
\usepackage{hyperref}
\usepackage{booktabs}
\usepackage{tabularx}
\usepackage{enumitem}
\usepackage{parskip}
\usepackage[most]{tcolorbox}
\usepackage{tikz}
\usetikzlibrary{arrows.meta,positioning}
\usepackage{colortbl}

\geometry{top=2.5cm, bottom=2.5cm, left=3cm, right=2.5cm}

\hypersetup{
  colorlinks=true,
  linkcolor=black,
  citecolor=blue!70!black,
  urlcolor=blue!70!black
}

\addbibresource{bibliography.bib}
\onehalfspacing

\begin{document}

% ---------------------------------------------------------------------------
% Titelseite
% ---------------------------------------------------------------------------
\begin{titlepage}
  \centering

  \large{HTWG Konstanz}

  \vfill

  {\LARGE\bfseries Ein lebender Wissensgraph für Coding Agents:\\[0.4em]
  Kontexttransparenz und strukturelle Context Assembly\\[0.4em]
  in der agentischen Softwareentwicklung\par}

  \vspace{1.2cm}

  {\large Exposé zur Masterarbeit}\\[0.4em]
  {\normalsize Betreuer-Ausschreibung: \textit{„Mehr Effizienz mit Coding Agents? Kontexttransparenz und Kontext-Engineering"}}

  \vspace{2cm}

  {\large Joschka Peeters}\\[0.5em]
  {\normalsize Matrikelnummer: 301750}\\[0.5em]
  {\normalsize \today}

  \vfill

  \textit{Betreuer: Prof. Dr. Rainer Müller}

\end{titlepage}

% ---------------------------------------------------------------------------
% Inhaltsverzeichnis
% ---------------------------------------------------------------------------
\tableofcontents
\newpage

% ---------------------------------------------------------------------------
% Sektionen
% ---------------------------------------------------------------------------
\section{Problemstellung und Motivation}

Coding Agents sind in der professionellen Softwareentwicklung angekommen.
Coding Agents wie SWE-agent übernehmen nicht länger nur Autovervollständigung,
sondern planen Aufgaben, rufen Werkzeuge auf, lesen Dateien
und koordinieren Änderungen über mehrere Module hinweg~\parencite{yang2024sweagent}.
Kommerzielle Werkzeuge wie Claude Code, Cursor und GitHub Copilot Workspace
setzen dieses Paradigma mittlerweile im Produktiveinsatz um.
An der Literatur lässt sich ein messbarer Produktivitätsgewinn belegen:
Entwicklerinnen und Entwickler schließen abgegrenzte Aufgaben mit AI-Unterstützung
bis zu 55,8\,\% schneller ab als ohne \parencite{peng2023copilot}.
Der SWE-bench-Benchmark zeigt, dass selbst State-of-the-Art-Modelle
an realen Repository-Aufgaben nahezu vollständig scheitern:
Das beste getestete Modell (Claude~2) löst lediglich 1,96\,\% der Issues~\parencite{jimenez2024swebench},
weil die Aufgaben Multi-File-Koordination, lange Kontexte und komplexes Reasoning erfordern.

Ein wesentliches Hindernis liegt dabei nicht allein im Modell,
sondern in der Kontextbereitstellung:
Was ein Coding Agent über eine Codebase weiß
(welche Dateien er gelesen hat, welche Abhängigkeiten er kennt,
welche Konventionen ihm vertraut sind)
bestimmt entscheidend, ob er eine Aufgabe korrekt löst
oder in produktionsfremden Mustern halluziniert~\parencite{li2026contextbench}.

\subsection*{Das Reaktivitätsproblem moderner Coding Agents}

Aktuelle Coding Agents explorieren Codebases reaktiv.
Das bedeutet, dass der Agent eigenständig Dateien öffnet und nach relevanten Funktionen sucht,
wenn er eine konkrete Aufgabe erhält.
Diese Strategie ist teuer und fehleranfällig.
\Textcite{li2026contextbench} dokumentieren, dass Coding Agents systematisch
mehr Kontext abrufen als sie tatsächlich nutzen,
und dabei vorher besuchte Dateien wiederholt aufsuchen
(ein Muster redundanter Exploration).
\Textcite{suri2026codescout} zeigen, dass Agents bei unklaren Aufgabenbeschreibungen
in eine ziellose Überexploration geraten.
Ein schlecht formulierter Bug-Report führt in ihren Experimenten
zu 21 Suchschritten ohne Ergebnis.
Mit einer strukturierten Aufgabenbeschreibung löst derselbe Agent
das gleiche Problem in sechs Schritten.

Es gibt bereits proaktive Ansätze, die Kontext vor der Agentausführung aufbereiten.
\textcite{xia2024agentless} lokalisiert Fehlerstellen in einem vorgelagerten Dreiphasen-Prozess.
\textcite{suri2026codescout} reichert die Aufgabenbeschreibung
durch eine Vorab-Exploration des Repositories an.
Beide Systeme analysieren die Codebase aufgabenbezogen
und zum Zeitpunkt der jeweiligen Anfrage.
Eine fortlaufende Synchronisation mit einer aktiv weiterentwickelten Codebase
ist nicht vorgesehen.

\subsection*{Das Transparenzproblem}

Neben dem Reaktivitätsproblem stellt sich die Frage der Inspizierbarkeit.
Wenn ein Coding Agent eine falsche Lösung produziert,
kann der Entwickler nicht nachvollziehen,
welcher Kontext zu welcher Entscheidung geführt hat.
Kontext ist eine Black Box.
Fehlt einer Funktion die Kenntnis ihrer Aufrufer?
Hat der Agent das zugehörige Schema nicht gelesen?
Ohne Antwort auf diese Fragen bleibt die Fehlersuche schwer nachvollziehbar.
Erste Ansätze zur Agenten-Transparenz, etwa AgentStepper~\parencite{agentstepper2026},
helfen Entwicklern, Ausführungstrajektorien nachzuvollziehen.
Die Inspizierbarkeit der Kontextauswahl selbst
findet dabei bislang kaum Beachtung.

\subsection*{Zielsetzung}

Ziel dieser Arbeit ist die Entwicklung und Evaluation eines Living Knowledge Graph
als Kontextinfrastruktur für Coding Agents.
Das System repräsentiert eine Codebase als strukturierten,
automatisch synchronisierten Graphen aus Dateien, Funktionen, Abhängigkeiten
und Testergebnissen.
Erhält das System eine Aufgabe, sucht es anhand des Living Knowledge Graph selbstständig die passenden Code-Abschnitte heraus
und gibt sie dem Agenten weiter.
Es erklärt dabei, warum jede Stelle ausgewählt wurde.
Das System lässt sich direkt in bestehende Werkzeuge wie Claude Code und Cursor einbinden.

Die Evaluation erfolgt an laufenden Hochschulprojekten,
um den Mehrwert im realen Entwicklungsalltag zu messen.
Diese Arbeit soll zeigen, dass Coding Agents durch den Living Knowledge Graph
weniger Zeit damit verbringen, sich in einer Codebase zu orientieren,
und weniger unnötige Dateien in ihren Kontext aufnehmen.
Außerdem sollen Entwickler besser nachvollziehen können,
warum ein Agent eine bestimmte Lösung vorschlägt.

\section{Stand der Forschung}

\subsection*{Agentische Ansätze für Software Engineering}

SWE-bench ist ein Datensatz aus 2\,294 realen GitHub-Issues mit zugehörigen
Codebase-Snapshots, Pull Requests und Test-Suites aus 12 Python-Repositories
\parencite{jimenez2024swebench}.
Der Agent erhält ein Issue und muss einen Patch generieren,
der die zugehörigen fail-to-pass-Tests besteht.
Die Aufgaben erfordern Änderungen in mehreren Dateien, komplexes Reasoning
und Interaktion mit Ausführungsumgebungen.
Klassische Sprachmodelle ohne Werkzeuge schaffen das nicht.

\textcite{yang2024sweagent} zeigen mit SWE-agent, dass ein spezialisiertes
Agent-Computer Interface (ACI, eine strukturierte Schnittstelle zwischen
Modell und Dateisystem) die Lösungsrate auf SWE-bench
von 3,8\,\% (bisherige RAG-Baseline) auf 12,5\,\% anhebt.
SWE-agent steht für den reaktiven Ansatz.
Der Agent entscheidet autonom, welche Dateien er liest,
welche Suchen er ausführt und welche Schritte er unternimmt.

Einen gegensätzlichen Ansatz verfolgt \textcite{xia2024agentless} mit Agentless.
Statt autonomer Exploration verwendet ein dreiphasiger, extern verwalteter Prozess
(Lokalisierung, Patchgenerierung, Testvalidierung) strukturierte Kontextzusammenstellung.
Auf SWE-bench Lite (einer 300-Aufgaben-Teilmenge des vollen SWE-bench)
erreicht Agentless 32\,\% bei nur 0,70\,USD pro Fix.
Das zeigt, dass prozedurale Kontextvorbereitung autonome Agentenexploration schlagen kann.


\subsection*{Wissensgraph-basierte Ansätze}

\textcite{ouyang2024repograph} stellen 2025 RepoGraph vor,
ein Plug-in-Modul zu bestehenden Software-Engineering-Agenten.
Es analysiert den Quellcode via AST (Abstract Syntax Tree,
ein automatisch erstellter Syntaxbaum)
und baut daraus einen Repository-Graphen aus Referenz- und Definitionsknoten
mit Aufruf- und Enthaltensrelationen.
Als Plug-in steigert RepoGraph die Lösungsrate von vier Baselines
(RAG, Agentless, AutoCodeRover, SWE-agent) auf SWE-bench Lite
um durchschnittlich 32,8\,\% relativ gegenüber der jeweiligen Baseline,
ohne die Kernsysteme zu verändern.

\textcite{yang2025kgcompass} entwickeln den graphbasierten Ansatz mit KGCompass weiter. Der Wissensgraph verknüpft nicht nur Code-Struktur,
sondern auch Issues und Pull Requests als Knoten im selben Graphen.
Dadurch lässt sich von einer natürlichsprachigen Bug-Beschreibung
direkt über Mehr-Hop-Pfade zu den relevanten Codestellen traversieren.
KGCompass erreicht 58,3\,\% auf SWE-bench Lite
und eine Fehlerlokalisierungsgenauigkeit von 56,0\,\%.
89,7\,\% der korrekt reparierten Patches erfordern dabei
Mehr-Hop-Verbindungen im Graphen.
Diese Codestellen sind im ursprünglichen Issue nicht direkt erwähnt,
sondern nur über mehrere Abhängigkeitsstufen erreichbar.

\subsection*{Proaktive versus reaktive Kontextbereitstellung}

Die Spannung zwischen proaktiver Vorbereitung und reaktiver Exploration
ist ein eigenständiges Forschungsthema.
\textcite{suri2026codescout} zeigen mit CodeScout,
dass proaktive Kontextanreicherung vor der Agentausführung
der reaktiven Eigenexploration überlegen ist.
Letztere neigt zu Über-Exploration und repetitiven Fix-Mustern,
während CodeScout durch strukturierte Vorab-Exploration präziseren Kontext liefert.

Ergänzend dokumentieren \textcite{li2026contextbench} mit ContextBench
auf 1\,136 Issue-Aufgaben aus 66 Repositories in acht Programmiersprachen:
Aktuelle Coding Agents priorisieren systematisch Recall über Präzision
(es wird mehr abgerufen als nötig),
und es bestehen erhebliche Lücken zwischen dem explorierten
und dem tatsächlich genutzten Kontext.
Komplexere Agent-Architekturen verbessern die Kontextabrufqualität dabei nur marginal.

Das von \textcite{liu2023lostmiddle} dokumentierte Lost-in-the-Middle-Phänomen
gibt diesem Befund eine modellseitige Erklärung.
Selbst als Long-Context-tauglich beworbene Modelle nutzen Informationen
in der Mitte langer Eingaben deutlich schlechter als an Anfang und Ende.
Proaktive, strukturell geordnete Kontextzusammenstellung ist daher
nicht nur eine Effizienzfrage, sondern eine Voraussetzung für zuverlässige Modellnutzung.

\subsection*{Statische Kontext-Dateien als Agenten-Baseline}

Eine weitere Herangehensweise untersucht den Wert manuell gepflegter Kontext-Dateien
(\texttt{CLAUDE.md}, \texttt{AGENTS.md}, \texttt{copilot-instructions.md})
für Coding Agents.
\textcite{gloaguen2026agentsmd} zeigen in einer systematischen Studie auf SWE-bench
(AGENTbench, 138 Aufgaben aus 12 Repositories),
dass vom Entwickler verfasste \texttt{AGENTS.md}-Dateien
die Lösungsrate um durchschnittlich 4\,\% verbessern.
LLM-generierte Dateien dagegen senken sie um 3\,\%.
Beide Varianten erhöhen die Inferenzkosten um mehr als 20\,\%.
\textcite{lulla2026agentsmd} ergänzen auf Basis von 124 Pull Requests aus 10 Repositories,
dass \texttt{AGENTS.md}-Dateien die mediane Bearbeitungszeit um 28,6\,\%
und den Output-Token-Verbrauch um 16,6\,\% reduzieren. 
Das bedeutet, dass Agenten mit besserem Kontext weniger Schritte benötigen, um eine Aufgabe zu lösen.

Diese Befunde belegen: Statische Entwickler-Dokumentation hat einen messbaren,
aber begrenzten Effekt auf Lösungsrate und Effizienz.
Die dargestellten Studien zu Kontext-Dateien vergleichen diese
ausschließlich mit einer Baseline ohne Kontext-Dateien.
Ein direkter Vergleich zwischen statischen Kontext-Dateien
und graphbasierter Kontextzusammenstellung findet sich
in der hier zitierten Literatur nicht.

\subsection*{Transparenz von Coding Agents}

Eine weitere offene Frage betrifft die Inspizierbarkeit agentischer Prozesse.
\textcite{agentstepper2026} stellen AgentStepper vor,
einen interaktiven Debugger für Coding Agents (SWE-agent, ExecutionAgent, RepairAgent),
der Entwicklern ermöglicht, Ausführungs-Trajektorien
(die Abfolge aller Schritte, die ein Agent bei einer Aufgabe durchführt)
schrittweise zu durchsuchen, Prompt-Verläufe einzusehen
und Agenten-Entscheidungen nachzuverfolgen.
In einer kontrollierten Studie mit 12 Teilnehmenden
stieg die Bug-Erkennungsrate von 17\,\% auf 60\,\%.
Diese Arbeit belegt, dass Einblick in die Ausführungsschritte
die Fähigkeit von Entwicklern zur Fehlerkorrektur messbar verbessert.

AgentStepper zeigt, welche Tool-Calls der Agent ausgeführt hat
und welche Ergebnisse sie lieferten.
Das im Rahmen dieser Arbeit zu entwickelnde System würde damit
auch über AgentStepper inspizierbar sein,
da seine Ergebnisse explizite Begründungen der Kontextauswahl enthalten.
Warum wurde genau diese Codestelle inkludiert?
Über welche Graph-Kante wurde sie erreicht?
Warum wurde eine andere Stelle ausgeschlossen?
KGCompass bettet Entscheidungspfade als Prompt-Kontext ein,
um die Patch-Generierung des LLM zu verbessern \parencite{yang2025kgcompass}.
Die Pfade dienen der Modellqualität, nicht der Entwicklerkontrolle.
Ob ein Entwickler die Kontextauswahl vor der Agentausführung bewerten
und korrigieren kann, wird nicht untersucht.


\subsection*{Integrationsinfrastruktur: MCP}

Im November 2024 hat Anthropic das Model Context Protocol (MCP)
als offenen Standard veröffentlicht \parencite{anthropic2024mcp}.
MCP definiert ein JSON-RPC-basiertes Protokoll zwischen einem Client (z.\,B. Claude Code)
und einem Server, der Kontext bereitstellt.
Ein Server kann drei Typen von Primitiven anbieten:
ausführbare Funktionen (Tools),
Datenquellen wie Dateiinhalte oder Datenbankrecords (Resources)
und wiederverwendbare Prompt-Templates auf dem Server (Prompts).
Für diese Arbeit ist MCP die natürliche Integrationsschicht:
Der Living-Knowledge-Graph-Server stellt seinen \texttt{get\_context}-Aufruf
als MCP-Tool bereit und ist damit ohne Änderung am Agent
in Claude Code, Cursor und kompatible Clients einbindbar.

\subsection*{Was diese Arbeit untersucht}

Die dargestellten Arbeiten belegen:
Graphbasierte Kontextzusammenstellung verbessert die Agent-Performance messbar.
Proaktive Kontextbereitstellung schlägt reaktive Exploration in Präzision und Effizienz.
Statische Entwickler-Dokumentation hat einen messbaren, aber begrenzten Effekt.
Diese Befunde basieren ausnahmslos auf Benchmark-Evaluationen auf SWE-bench.
SWE-bench-Aufgaben basieren auf eingefrorenen Codebase-Snapshots mit vorgefertigten Testsuiten.
Der Mehrwert eines kontinuierlich synchronisierten Wissensgraphen lässt sich dort nicht messen.
Laufende Hochschulprojekte werden aktiv weiterentwickelt, haben keine vollständige Testabdeckung
und involvieren menschliche Entwickler mit unterschiedlichem Erfahrungsstand.
Ob ein Living Knowledge Graph auch unter diesen Bedingungen einen messbaren Mehrwert liefert,
ist bislang nicht untersucht.

Daraus ergeben sich drei Fragen:

Wie lässt sich ein Wissensgraph dauerhaft mit einer aktiv weiterentwickelten
Codebase synchron halten?
Bestehende Systeme analysieren das Repository zum Aufgabenzeitpunkt;
eine inkrementelle, commit-getriebene Aktualisierung fehlt.

Verbessert ein Living Knowledge Graph die Arbeitseffizienz von Studierenden
an realen Softwareprojekten messbar,
verglichen mit reaktiver Exploration und statischer Dokumentation?

Ermöglicht die explizite Graph-Pfad-Begründung im MCP-Response Entwicklern,
Kontextentscheidungen des Agenten vor der Ausführung zu bewerten
und bei Bedarf zu korrigieren?

Die vorliegende Arbeit beantwortet diese Fragen mit einem Living Knowledge Graph,
der inkrementell mit der Codebase synchronisiert wird
und als MCP-Server in Claude Code, Cursor und kompatible Umgebungen einbindbar ist.
Jeder gelieferte Kontext enthält einen expliziten Entscheidungspfad,
sodass Entwickler die Kontextauswahl nachvollziehen und bei Bedarf eingreifen können.

\section{Forschungsfragen und Hypothesen}

\subsection*{Forschungsfragen}

\begin{description}
  \item[HF1] Verbessert proaktiv graph-assemblierter Kontext die Effizienz von Coding Agents (gemessen an Token-Verbrauch, Iterationsanzahl und Erstkorrektheit, d.\,h. Rechenkosten, Agentenschritten und Erstlösungsqualität) gegenüber reaktiver Agent-Eigenexploration (B1) und gegenüber statischer Entwickler-Dokumentation (B2)?

  \item[HF2] Welche Graph-Traversal-Strategien (strukturell-deklarativ, historisch, testbasiert) liefern für welche Aufgabentypen (isolierte Bugfixes, Multi-File-Refactorings, neue Feature-Implementierungen) den präzisesten Kontext?

  \item[HF3] Erhöht die transparente Visualisierung der Kontextauswahl die Fähigkeit von Entwicklern, Fehler im Agent-Output zu erkennen und gezielt zu korrigieren?
\end{description}

\subsection*{Hypothesen}

\begin{description}
  \item[H1] Graph-basiert proaktiv assemblierter Kontext reduziert Token-Verbrauch und Iterationsanzahl gegenüber reaktiver Exploration und gegenüber statischer Entwickler-Dokumentation. Diese Hypothese stützt sich auf drei komplementäre Befunde: \textcite{li2026contextbench} dokumentieren, dass reaktive Agents erheblich mehr Kontext abrufen als sie nutzen und vorher besuchte Dateien wiederholt aufsuchen. \textcite{xia2024agentless} zeigen, dass prozedurale Kontextvorbereitung reaktive Agentenexploration auf SWE-bench Lite mit 32\,\% und nur 0,70\,USD pro Fix übertrifft. \textcite{gloaguen2026agentsmd} zeigen, dass statische Kontext-Dateien zwar die Lösungsrate um 4\,\% verbessern, die Inferenzkosten aber um mehr als 20\,\% erhöhen, was nahelegt, dass aufgabenspezifisch gefilterte Kontextzusammenstellung effizienter sein kann. Ein strukturell aufgebauter Graph, der nur den direkt abhängigen Teilgraphen ausliefert, sollte dieses Effizienzpotenzial ausschöpfen.

  \item[H2] Multi-Hop-Graphtraversierung über Aufruf- und Abhängigkeitskanten liefert für Multi-File-Aufgaben präziseren Kontext als flaches Dateisuchen. Diese Hypothese basiert auf dem zentralen KGCompass-Befund: 89,7\,\% der korrekt reparierten Patches erfordern Mehr-Hop-Verbindungen im Graphen und betreffen Codestellen, die reine LLMs ohne Graphtraversierung systematisch übersehen \parencite{yang2025kgcompass}. Der relative Leistungsgewinn von RepoGraph über bestehende Systeme (+32,8\,\%) bestätigt den Mehrwert struktureller Repräsentation \parencite{ouyang2024repograph}.

  \item[H3] Entwickler, denen die \emph{Retrieval-Transparenz} (also welche Graph-Knoten warum in den Kontext aufgenommen wurden) als inspizierbarer Entscheidungspfad präsentiert wird, erkennen mehr Agent-Fehler und benötigen weniger Korrekturiterationen als Entwickler ohne diese Information. Die Hypothese ist explorativ. Verwandte Arbeiten zeigen, dass Transparenz über Ausführungs-Trajektorien die Fehlererkennungsrate messbar verbessert \parencite{agentstepper2026}; die vorgelagerte Frage der Retrieval-Transparenz ist bislang empirisch unerforscht. Ergänzend zeigt \textcite{liu2023lostmiddle}, dass fehlerhafte Kontextpositionierung zu Modellfehlern führt, was nahelegt, dass sichtbarer Kontext Entwicklern gezielteren Eingriff ermöglicht.
\end{description}

\section{Das Artefakt: Living Knowledge Graph}

Das zentrale Artefakt dieser Arbeit ist ein System, das eine Codebase als strukturierten Wissensgraphen repräsentiert und diesen als proaktive Kontextinfrastruktur für Coding Agents bereitstellt. \emph{Living} bezeichnet dabei die Eigenschaft, dass der Graph nach jedem Commit automatisch aktualisiert wird und damit dauerhaft den aktuellen Codestand abbildet (siehe Schicht~2). Das System gliedert sich in drei Schichten; Abbildung~\ref{fig:architecture} gibt einen Überblick.

\begin{figure}[htbp]
\centering
\begin{tikzpicture}[
  node distance=0.65cm,
  mybox/.style={
    draw, rounded corners=3pt,
    minimum width=12cm, align=center,
    inner sep=9pt, font=\small
  },
  arr/.style={-{Stealth[length=5pt,width=4pt]}, thick, gray!55},
  lbl/.style={font=\scriptsize\itshape, text=gray!65}
]

\node[mybox, fill=blue!8] (codebase) {%
  \textbf{Quellcode-Repositories}\quad
  \texttt{.ts}\;/\;\texttt{.tsx}\;/\;\texttt{.py}\;/\;\texttt{.js}};

\node[mybox, fill=orange!10, below=of codebase] (s1) {%
  \textbf{Schicht 1: Ingestion (tree-sitter)}\\[3pt]
  Knoten:\enspace\texttt{File},\enspace\texttt{Function},\enspace\texttt{Class}%
  \qquad
  Kanten:\enspace\texttt{IMPORTS},\enspace\texttt{CALLS},\enspace\texttt{EXTENDS},\enspace\texttt{DEFINED\_IN}};

\node[mybox, fill=orange!14, below=of s1] (s2) {%
  \textbf{Schicht 2: Living Graph (Neo4j)}\\[3pt]
  Git post-commit Hook\enspace$\cdot$\enspace
  inkrementelle Synchronisation\enspace$\cdot$\enspace
  Git-History-Metriken};

\node[mybox, fill=orange!19, below=of s2] (s3) {%
  \textbf{Schicht 3: Context Assembly}\\[3pt]
  Entity Extraction\;$\to$\;Seed-Lookup\;$\to$\;
  Graph-Traversal\;$\to$\;Ranking\;$\to$\;Entscheidungspfad};

\node[mybox, fill=teal!10, below=of s3] (mcp) {%
  \textbf{MCP-Server}\quad
  \texttt{get\_context(task)}\enspace
  \texttt{inspect\_graph(node\_id)}\enspace
  \texttt{update\_graph()}};

\node[mybox, fill=green!8, below=of mcp] (agent) {%
  \textbf{Coding Agent}\quad
  Claude Code\;$\cdot$\;Cursor\;$\cdot$\;weitere MCP-kompatible Clients};

\draw[arr] (codebase) -- node[lbl, right=5pt]{tree-sitter Parser} (s1);
\draw[arr] (s1)       -- node[lbl, right=5pt]{Graph aufbauen} (s2);
\draw[arr] (s2)       -- node[lbl, right=5pt]{Graphabfrage via Cypher} (s3);
\draw[arr] (s3)       -- node[lbl, right=5pt]{MCP JSON-RPC} (mcp);
\draw[arr] (mcp)      -- node[lbl, right=5pt]{Kontext-Block und Entscheidungspfad} (agent);

\end{tikzpicture}
\caption{Systemarchitektur des \emph{Living Knowledge Graph}. Der Coding Agent erhält
über den MCP-Server proaktiv assemblierten, aufgabenspezifischen Kontext, ohne die
Codebase reaktiv zu explorieren.}
\label{fig:architecture}
\end{figure}


\subsection*{Schicht 1: Ingestion und Graph-Aufbau}

Ein auf \emph{tree-sitter} basierender Parser verarbeitet Python-, JavaScript- und TypeScript-Quellcode und extrahiert daraus einen strukturierten Graphen in Neo4j:

\begin{tabularx}{\linewidth}{@{}lX@{}}
  \toprule
  \textbf{Knotentyp} & \textbf{Bedeutung} \\
  \midrule
  \texttt{File}      & Quellcode-Datei mit Pfad und Sprache \\
  \texttt{Function}  & Funktion oder Methode mit Signatur und Zeilenbereich \\
  \texttt{Class}     & Klasse oder Interface \\
  \texttt{TestCase}  & Testfunktion mit Referenz zum getesteten Symbol \\
  \texttt{Schema}    & Daten- oder API-Schema (z.\,B. Sanity- oder Pydantic-Schema) \\
  \bottomrule
\end{tabularx}

\vspace{0.5em}

\begin{tabularx}{\linewidth}{@{}lX@{}}
  \toprule
  \textbf{Kantentyp} & \textbf{Bedeutung} \\
  \midrule
  \texttt{IMPORTS}   & Datei importiert ein Symbol aus einer anderen Datei \\
  \texttt{CALLS}     & Funktion ruft eine andere Funktion auf \\
  \texttt{EXTENDS}    & Klasse erbt von einer anderen Klasse \\
  \texttt{DEFINED\_IN} & Funktion oder Klasse ist in einer bestimmten Datei definiert \\
  \texttt{TESTS}     & Testfunktion testet ein bestimmtes Symbol \\
  \texttt{DOCUMENTS} & Docstring oder Kommentar dokumentiert ein Symbol \\
  \bottomrule
\end{tabularx}

\vspace{0.5em}

Ergänzend werden Git-History-Metriken als Knotenattribute gespeichert: Änderungshäufigkeit und letzte Commit-Nachricht. Diese ermöglichen risikogewichtetes Ranking: Häufig modifizierte Dateien werden bei gleicher Hop-Distanz bevorzugt.

\subsection*{Schicht 2: Living Graph mit inkrementeller Git-Synchronisation}

Ein Git post-commit Hook analysiert nach jedem Commit, welche Dateien sich verändert haben. Nur die betroffenen Knoten und ihre direkten Kanten werden neu geparst und aktualisiert; ein vollständiger Rebuild entfällt. Damit bleibt der Graph dauerhaft synchron mit dem Codestand, ohne dass manuelle Wartung nötig ist. Dies ist die \emph{Living}-Eigenschaft des Systems: Der Wissensgraph ist kein statisches Dokument, sondern ein kontinuierlich aktualisiertes Abbild der Codebase.

\subsection*{Schicht 3: Proaktive Context Assembly}

Für eine eingehende Aufgabenbeschreibung läuft folgende Pipeline:

\begin{enumerate}
  \item \textbf{Entity Extraction:} Aus dem Aufgabentext werden genannte Funktionsnamen, Dateinamen und Schlüsselbegriffe extrahiert (via kleinem LLM-Call oder Regex-Matching).
  \item \textbf{Seed-Node-Identifikation:} Die extrahierten Entitäten werden im Graphen als Startknoten lokalisiert.
  \item \textbf{Graph-Traversal:} Von den Startknoten aus wird der Graph in konfigurierbarer Hop-Tiefe traversiert. Unterschiedliche Traversal-Strategien sind wählbar: Strukturell (Aufruf- und Importkanten), historisch (häufig co-modifizierte Knoten) und testbasiert (zugehörige Testknoten).
  \item \textbf{Ranking und Filterung:} Knoten werden nach Relevanz geordnet (Hop-Distanz, Testabdeckung, Risikogewicht) und auf ein konfigurierbares Tokenbudget begrenzt.
  \item \textbf{Serialisierung:} Der Teilgraph wird als strukturierter Kontext-Block (Markdown oder JSON) ausgegeben, inklusive expliziter Angaben, \emph{warum} jeder Knoten aufgenommen wurde.
\end{enumerate}

\begin{figure}[htbp]
\centering
\begin{tikzpicture}[
  node distance=0.45cm,
  pbox/.style={
    draw, rounded corners=3pt,
    minimum width=2.15cm, minimum height=1.1cm,
    align=center, inner sep=4pt,
    font=\small\bfseries, fill=orange!12
  },
  exlbl/.style={font=\scriptsize, text=gray!70, align=center},
  arr/.style={-{Stealth[length=5pt,width=4pt]}, thick, gray!55}
]

%% Pipeline boxes
\node[pbox] (e1) {Entity\\Extraction};
\node[pbox, right=of e1] (e2) {Seed-\\Lookup};
\node[pbox, right=of e2] (e3) {Graph-\\Traversal};
\node[pbox, right=of e3] (e4) {Ranking\\\& Filter};
\node[pbox, right=of e4] (e5) {Seriali-\\sierung};

%% Arrows between boxes
\draw[arr] (e1) -- (e2);
\draw[arr] (e2) -- (e3);
\draw[arr] (e3) -- (e4);
\draw[arr] (e4) -- (e5);

%% Input / output labels above
\node[exlbl, above=5pt of e1] {\textit{Aufgabe}};
\node[exlbl, above=5pt of e5] {\textit{Kontext-Block}};

%% Example row below
\node[exlbl, below=8pt of e1] {\glqq{}fix cart\\checkout\grqq{}};
\node[exlbl, below=8pt of e2] {\{checkout,\\cart, Cart\}};
\node[exlbl, below=8pt of e3] {32 Knoten,\\2 Hops};
\node[exlbl, below=8pt of e4] {Top 15,\\sortiert};
\node[exlbl, below=8pt of e5] {Markdown\\+~Pfad-Audit};

%% Dotted separator line
\draw[dotted, gray!50]
  ([xshift=-0.3cm, yshift=-0.35cm]e1.south west) --
  ([xshift= 0.3cm, yshift=-0.35cm]e5.south east);

\end{tikzpicture}
\caption{Fünfstufige Context-Assembly-Pipeline (\texttt{get\_context}). Die untere Zeile
zeigt ein Beispiel für die Aufgabe \emph{,,fix cart checkout``}.}
\label{fig:pipeline}
\end{figure}


\subsection*{Transparenz als First-Class-Feature}

Jeder ausgelieferte Kontext-Block enthält einen \emph{Entscheidungspfad}: Eine maschinenlesbare Begründung, über welche Kanten jeder Knoten erreicht wurde (\texttt{checkout() aufgenommen via CALLS-Kante von cart.ts, Zeile 42}). Entwickler können diesen Teilgraphen vor Agentausführung inspizieren und gezielt erweitern oder beschneiden. Nach Abschluss der Aufgabe dient der Entscheidungspfad als Audit-Trail.

\subsection*{Integration via MCP}

Das System wird als \emph{MCP-Server} implementiert und stellt drei Tools bereit:

\begin{description}
  \item[\texttt{get\_context(task)}] Assembliert und gibt den Kontext-Block für eine Aufgabenbeschreibung zurück.
  \item[\texttt{inspect\_graph(node\_id)}] Gibt den direkten Nachbarschaft-Subgraphen eines Knotens zurück.
  \item[\texttt{update\_graph()}] Triggert manuell eine Graph-Aktualisierung (außerhalb des automatischen Git-Hooks).
\end{description}

Claude Code, Cursor und andere MCP-kompatible Clients können diese Tools direkt aufrufen. Damit ist das System ohne Änderung am Agent selbst nutzbar.

\subsection*{Stand der Vorarbeiten}

\begin{tcolorbox}[colback=blue!4!white, colframe=blue!40!black, title=Proof-of-Concept bereits implementiert, fonttitle=\bfseries]
Im Rahmen der Einarbeitungsphase wurde ein funktionsfähiger Proof-of-Concept für die TypeScript-Komponente entwickelt und auf dem eigenen WebshopTemplate-Projekt validiert:

\begin{itemize}[noitemsep]
  \item \textbf{Parser (M2):} tree-sitter-Parser für TypeScript/TSX extrahiert Datei-, Funktions- und Klassenknoten sowie IMPORTS-, CALLS-, EXTENDS- und DEFINED\_IN-Kanten.
  \item \textbf{Graph (M3):} Neo4j-Instanz mit 59 Dateiknoten, 64 Funktionsknoten und 25 aufgelösten CALLS-Kanten; Git-History-Metriken (\texttt{commitCount}, \texttt{lastCommitMessage}) als Knotenattribute; Performance-Indexe.
  \item \textbf{Context Assembly (M5):} Fünfstufige Pipeline (\texttt{get\_context}) mit Entity Extraction, Seed-Node-Identifikation, konfigurierbarer Graphtraversierung, Relevanz-Ranking und Markdown-Serialisierung inklusive Entscheidungspfad.
  \item \textbf{MCP-Server (M6):} Drei Tools (\texttt{get\_context}, \texttt{inspect\_graph}, \texttt{update\_graph}) als MCP-Server implementiert und in Claude Code registriert.
\end{itemize}

Die Living-Eigenschaft (automatischer Git-Hook, M4) und die Python-Komponente sind in der eigentlichen Masterarbeit zu entwickeln.
\end{tcolorbox}

\section{Methodisches Vorgehen}

\subsection*{Forschungsrahmen: Design Science Research}

Die Arbeit folgt dem \emph{Design Science Research} (DSR)-Paradigma nach \textcite{hevner2004dsr}. DSR ist für informatiknahe Masterarbeiten geeignet, die ein nützliches Artefakt entwickeln und dessen Wirksamkeit empirisch nachweisen. Der Zyklus umfasst drei Phasen: \emph{Relevance} (das Problem ist real und praxisrelevant), \emph{Design} (das Artefakt wird iterativ entwickelt und verfeinert) und \emph{Rigor} (die Evaluation ist methodisch fundiert). Für diese Arbeit bedeutet das: Das System wird nicht einmalig gebaut und dann evaluiert, sondern über mehrere Build-Evaluate-Zyklen iteriert.

\subsection*{Evaluation: Drei Versuchsbedingungen}

Die Evaluation vergleicht drei Bedingungen:

\begin{description}
  \item[B1: Reaktive Exploration] Der Coding Agent erhält nur die Aufgabenbeschreibung und exploriert die Codebase selbständig (Baseline, entspricht dem aktuellen Stand der Praxis).
  \item[B2: Statischer Kontext] Der Agent erhält zusätzlich eine statisch gepflegte \texttt{CLAUDE.md} mit Architekturhinweisen und Konventionen. Diese Bedingung entspricht dem heutigen Best-Practice und ist empirisch fundiert: Entwickler-geschriebene Kontext-Dateien verbessern die Lösungsrate um durchschnittlich 4\,\%, reduzieren aber die Inferenzeffizienz \parencite{gloaguen2026agentsmd,lulla2026agentsmd}.
  \item[B3: Living Knowledge Graph] Der Agent erhält den durch das System assemblierten, aufgabenspezifischen Kontext-Block inklusive Entscheidungspfad (das neue Artefakt).
\end{description}

\begin{figure}[htbp]
\centering
\renewcommand{\arraystretch}{2.2}
\begin{tabular}{|>{\centering\arraybackslash}m{2.8cm}%
                |>{\centering\arraybackslash}m{3.2cm}%
                |>{\centering\arraybackslash}m{3.2cm}%
                |>{\centering\arraybackslash}m{3.2cm}|}
\hline
\rowcolor{gray!22}
\textbf{Bedingung}
  & \textbf{T1: Bugfixes}
  & \textbf{T2: Refactorings}
  & \textbf{T3: Features} \\
\hline
\cellcolor{gray!12}{\small\textbf{B1\newline Reaktive\newline Exploration}}
  & \cellcolor{red!5}{\small Baseline}
  & \cellcolor{red!5}{\small Baseline}
  & \cellcolor{red!5}{\small Baseline} \\
\hline
\cellcolor{gray!12}{\small\textbf{B2\newline Statisch\newline (CLAUDE.md)}}
  & \cellcolor{yellow!18}{\small HF1}
  & \cellcolor{yellow!18}{\small HF1}
  & \cellcolor{yellow!18}{\small HF1} \\
\hline
\cellcolor{gray!12}{\small\textbf{B3\newline Living\newline Knowledge Graph}}
  & \cellcolor{green!10}{\small HF1 $\cdot$ HF2}
  & \cellcolor{green!10}{\small HF1 $\cdot$ HF2}
  & \cellcolor{green!16}{\small HF1 $\cdot$ HF2 $\cdot$ \textbf{HF3}} \\
\hline
\end{tabular}

\smallskip
{\footnotesize\textit{Metriken je Zelle: Token-Verbrauch, Iterationsanzahl, Erstkorrektheit,
Kontextpräzision. HF3 (Retrieval-Transparenz) wird über alle B3-Zellen erhoben; T3 ist der
primäre Messkontext.}}

\caption{Evaluationsdesign: drei Versuchsbedingungen (B1 bis B3) gekreuzt mit drei
Aufgabentypen (T1 bis T3). Jede Zelle steht für einen Messdurchlauf.}
\label{fig:evaluation}
\end{figure}


\subsection*{Versuchsumgebungen}

Die Evaluation wird auf zwei bis drei laufenden Hochschulprojekten des Betreuers durchgeführt, also aktiv weiterentwickelten Python- und TypeScript-Projekten mit gewachsener Codebasis. Diese Umgebungen bieten realistische, nicht konstruierte Aufgaben und ermöglichen die Beteiligung von Studierenden als Probanden. Während industrielle Studien wie \textcite{hula2024} Coding Agents auf realen Projekten evaluieren und dabei Durchsatz und Akzeptanzraten messen, steht bei dieser Arbeit die Fähigkeit der Entwicklenden im Vordergrund, die Kontextauswahl des Agenten zu verstehen und gezielt zu korrigieren, eine Dimension, die in bestehenden Feldevaluationen bislang kaum erfasst wird.

Pro Projekt werden Aufgaben dreier Typen bearbeitet:
\begin{description}
  \item[T1: Isolierte Bugfixes] Lokal begrenzte Korrekturen.
  \item[T2: Multi-File-Refactorings] Strukturelle Änderungen mit dateiübergreifenden Auswirkungen.
  \item[T3: Feature-Entwicklungen] Neue Funktionalität mit Architektur- und Konventionsbedarf.
\end{description}

\subsection*{Metriken}

\begin{tabularx}{\linewidth}{@{}lXl@{}}
  \toprule
  \textbf{Dimension} & \textbf{Operationalisierung} & \textbf{Erhebung} \\
  \midrule
  Effizienz       & Iterationsanzahl bis zur akzeptierten Lösung & Logging \\
  Token-Verbrauch & Input- und Output-Tokens pro Aufgabe & API-Logs \\
  Erstkorrektheit & Anteil Aufgaben, die im ersten Durchlauf alle Tests bestehen & Test-Suite \\
  Kontextpräzision & Anteil relevanter zu insgesamt ausgelieferter Kontext-Tokens & Annotierung \\
  Transparenz & Entwickler-Bewertung der Verständlichkeit des Entscheidungspfads & Kurzumfrage \\
  \bottomrule
\end{tabularx}

\subsection*{Validität und Abgrenzung}

Die interne Validität wird durch reproduzierbare Aufgaben, fixierte Modellversionen und mehrfache Wiederholungen je Bedingung gesichert. Die externe Validität ergibt sich aus der Nutzung realer, aktiv betreuter Hochschulprojekte statt konstruierter Benchmarks. Eine Einschränkung ist die Beschränkung auf die Sprachen Python, JavaScript und TypeScript sowie auf MCP-kompatible Coding Agents.

\section{Vorläufige Gliederung der Arbeit}

\begin{description}
  \item[1. Einleitung] Motivation, Problemstellung, Zielsetzung und Aufbau der Arbeit.

  \item[2. Grundlagen] Große Sprachmodelle für Code; agentische Architekturen (ReAct, Werkzeugnutzung); Wissensgraphen und Graphdatenbanken; Kontext- und Prompt-Engineering; Model Context Protocol.

  \item[3. Stand der Forschung] Agentische Ansätze für Software Engineering (SWE-agent, Agentless); graphbasierte Kontextrepräsentation (RepoGraph, KGCompass); proaktive vs.\ reaktive Kontextbereitstellung (CodeScout, ContextBench); statische Kontext-Dateien als Agenten-Baseline (AGENTS.md-Literatur); Ableitung der Forschungslücke.

  \item[4. Systemarchitektur] Detaillierte Beschreibung der drei Schichten: Ingestion (tree-sitter, Parser-Design), Living Graph (Neo4j-Schema, Git-Synchronisation) und Context Assembly (Traversal-Strategien, Ranking, Serialisierung); MCP-Integration; Transparenzkonzept.

  \item[5. Implementierung] Technische Umsetzung in Python und TypeScript; Deployment auf den Hochschulprojekten; Logging-Pipeline; Aufgabenkatalog.

  \item[6. Ergebnisse] Quantitative Befunde (B1 vs.\ B2 vs.\ B3 nach Aufgabentyp und Metrik); qualitative Beobachtungen aus Agent-Logs; Transparenz-Feedback der Entwickler.

  \item[7. Diskussion] Interpretation der Ergebnisse; Rückbezug auf Hypothesen; Traversal-Strategien im Vergleich; Grenzen des Ansatzes; praktische Implikationen für Kontext-Engineering.

  \item[8. Fazit und Ausblick] Zusammenfassung der Beiträge; offene Forschungsfragen; mögliche Erweiterungen (weitere Sprachen, automatische Traversal-Parametrisierung, LLM-gestützte Graphanreicherung).

  \item[Anhang] Aufgabenkatalog; Graph-Schemata; Prompt-Templates; auszugsweise Roh-Logs; ergänzende Tabellen.
\end{description}

\section{Zeitplan}

Die Bearbeitung ist auf 24 Wochen angelegt. Die ersten sechs Wochen sind Grundlagen und Systemkern gewidmet, sodass ab Woche 7 ein laufendes System existiert, das iterativ auf den Hochschulprojekten verfeinert wird.

\begin{tabularx}{\linewidth}{@{}lXl@{}}
  \toprule
  \textbf{Wochen} & \textbf{Phase / Aktivität} & \textbf{Meilenstein} \\
  \midrule
  W01--W02 & Literaturvertiefung; Absprache Hochschulprojekte; DSR-Design finalisieren & M1: Design final \\
  W03--W04 & tree-sitter Parser für Python + TypeScript/JavaScript; Neo4j-Schema & M2: Parser läuft\textsuperscript{*} \\
  W05--W06 & Graph-Builder (Ingestion Layer); Erstbefüllung eines Hochschulprojekts & M3: Graph befüllt\textsuperscript{*} \\
  W07--W08 & Git post-commit Hook; inkrementelle Synchronisation (Living Layer) & M4: Living Graph aktiv \\
  W09--W10 & Context Assembly: Traversal-Strategien, Ranking, Serialisierung & M5: Assembly läuft\textsuperscript{*} \\
  W11--W12 & MCP-Server; Integration in Claude Code; erster End-to-End-Test & M6: MCP integriert\textsuperscript{*} \\
  W13--W15 & Transparenz-Features: Entscheidungspfade, Subgraph-Visualisierung & M7: Transparenz fertig \\
  W16--W18 & Evaluation auf allen Hochschulprojekten; Datenerhebung B1 vs.\ B2 vs.\ B3 & M8: Daten erhoben \\
  W19--W21 & Datenanalyse (quantitativ + qualitativ); Iteration am System & M9: Ergebnisse konsolidiert \\
  W22--W24 & Verschriftlichung; Review; Korrekturen; Abgabe & M10: Abgabe \\
  \bottomrule
\end{tabularx}

\noindent\textsuperscript{*} Im Rahmen der Einarbeitungsphase für TypeScript/WebshopTemplate bereits abgeschlossen (siehe Abschnitt~4).

\subsection*{Risikobetrachtung}

\begin{description}
  \item[R1: Parser-Abdeckungslücken] tree-sitter-Grammatiken decken nicht alle Python- und TypeScript-Konstrukte vollständig ab. \emph{Gegenmaßnahme:} Frühzeitiger Test auf den Ziel-Projekten in W03--W04; Fallback auf vereinfachtes Kanten-Set.

  \item[R2: Zugang zu Hochschulprojekten] Nicht alle Projekte sind sofort zugänglich. \emph{Gegenmaßnahme:} Frühzeitige Abstimmung in W01; paralleles Arbeiten mit dem eigenen WebshopTemplate als Ersatz-Lab.

  \item[R3: Neo4j-Performance bei großen Graphen] Bei Projekten mit mehr als 50\,000 Knoten können Traversal-Anfragen langsam werden. \emph{Gegenmaßnahme:} Indexierung kritischer Felder ab W06; Subgraph-Caching für häufige Queries.

  \item[R4: Modell-Updates] Anbieterseitige Änderungen an Claude Code oder der API können Evaluationsergebnisse beeinflussen. \emph{Gegenmaßnahme:} Modellversion und API-Version für die gesamte Erhebungsphase fixieren und dokumentieren.

  \item[R5: Token- und API-Kosten] Wiederholte agentische Läufe für B1--B3 können das Budget überschreiten. \emph{Gegenmaßnahme:} Vorabkalkulation pro Bedingung; Priorisierung der aussagekräftigsten Aufgaben-Typ-Kombinationen.
\end{description}

\newpage
\section{Einarbeitungshinweise}

Dieser Abschnitt dokumentiert die Einarbeitungsphase (W01--W02), deren praktische Ergebnisse als Proof-of-Concept in Abschnitt~4 beschrieben sind. Die Ressourcen sind nach Priorität geordnet.

% -----------------------------------------------------------------------
\subsection*{1. tree-sitter: Code-Parser}

tree-sitter ist der technische Kern der Ingestion-Schicht. Priorität: \textbf{hoch}.

\begin{description}
  \item[Einstieg] Offizielle Dokumentation: \url{https://tree-sitter.github.io/tree-sitter/using-parsers}. Den Abschnitt \emph{„Using Parsers"} und \emph{„The Node API"} vollständig lesen.

  \item[Python-Bindings] Paket \texttt{tree-sitter} + \texttt{tree-sitter-python} und \texttt{tree-sitter-typescript} installieren. Ein kleines Skript schreiben, das eine Python-Datei parst und alle Funktionsnamen + Signaturen ausgibt; das ist der Kern-Use-Case für diese Arbeit.

  \item[Übung] Aus einem echten Python-Projekt (z.\,B. dem WebshopTemplate) alle \texttt{import}-Statements extrahieren und als Kantenliste ausgeben: \texttt{datei\_a.py IMPORTS symbol\_x FROM datei\_b.py}.

  \item[Wichtig zu verstehen] Der Unterschied zwischen einem \emph{Syntax-Tree} (was tree-sitter liefert) und einem \emph{Scope-Tree} (was für Aufrufkanten nötig ist). Für einfache Aufrufkanten reicht tree-sitter; für präzises Cross-File-Calling ist zusätzliche Auflösung nötig.
\end{description}

% -----------------------------------------------------------------------
\subsection*{2. Neo4j: Graphdatenbank}

Neo4j ist die Persistenzschicht. Priorität: \textbf{hoch}.

\begin{description}
  \item[Kurs] Neo4j GraphAcademy (kostenlos): \url{https://graphacademy.neo4j.com}. Dort die Kurse \emph{„Neo4j Fundamentals"} und \emph{„Cypher Fundamentals"} absolvieren. Zusammen ca.\ 4 Stunden.

  \item[Lokaler Start] Neo4j Desktop herunterladen oder AuraDB Free Tier (Cloud, kein Setup): \url{https://neo4j.com/cloud/platform/aura-graph-database/}. AuraDB empfiehlt sich für den Einstieg.

  \item[Python-Driver] \texttt{neo4j} Python-Paket. Eine kleine Pipeline schreiben, die den tree-sitter-Output in Neo4j-Knoten und -Kanten überführt.

  \item[Wichtige Cypher-Abfragen üben:]
    \begin{itemize}
      \item Alle Funktionen die eine bestimmte Funktion aufrufen (eingehende CALLS-Kanten)
      \item n-Hop-Nachbarschaft eines Knotens (\texttt{MATCH path = (n)-[*1..3]->(m)})
      \item Kürzester Pfad zwischen zwei Knoten
    \end{itemize}
\end{description}

% -----------------------------------------------------------------------
\subsection*{3. Pflichtlektüre: Kernpaper}

Diese Paper sollten vor dem Beginn der Implementierung gelesen werden, in dieser Reihenfolge:

\begin{enumerate}
  \item \textbf{KGCompass} \parencite{yang2025kgcompass}: Direkt verwandter Ansatz; Graph-Schema und Multi-Hop-Traversal genau analysieren. Besonders: Abschnitte zu Graph-Konstruktion und Fault-Localization.

  \item \textbf{RepoGraph} \parencite{ouyang2024repograph}: Zeigt, wie ein Code-Graph als Plug-in-Modul funktioniert; Ego-Graph-Konzept ist relevant für die Context-Assembly-Schicht.

  \item \textbf{SWE-agent} \parencite{yang2024sweagent}: Zeigt wie reaktive Agent-Computer-Interfaces funktionieren; Grundlage für Bedingung B1 in der Evaluation.

  \item \textbf{Agentless} \parencite{xia2024agentless}: Zeigt, dass prozedurale Kontextvorbereitung reaktive Exploration schlagen kann; wichtig für die Argumentation in HF1.

  \item \textbf{ContextBench} \parencite{li2026contextbench}: Empirische Grundlage für das Over-Retrieval-Problem; Benchmark-Design studieren für eigene Evaluation.
\end{enumerate}

% -----------------------------------------------------------------------
\subsection*{4. Model Context Protocol (MCP)}

MCP ist die Integrationsschicht. Priorität: \textbf{mittel} (nach W10).

\begin{description}
  \item[Dokumentation] \url{https://modelcontextprotocol.io}: Abschnitte \emph{„Introduction"}, \emph{„Architecture"} und \emph{„Building Servers"} lesen.

  \item[Python-SDK] \url{https://github.com/modelcontextprotocol/python-sdk}: Das offizielle Tutorial \emph{„Building a simple MCP server"} vollständig durcharbeiten. Ein MCP-Server der ein einziges Tool (\texttt{hello\_world}) bereitstellt, in unter einer Stunde gebaut.

  \item[Integration in Claude Code] In Claude Code: \texttt{/mcp} Kommando zum Registrieren eigener Server. Der eigene MCP-Server kann damit direkt in der Entwicklungsumgebung getestet werden, ohne eine separate App.
\end{description}

% -----------------------------------------------------------------------
\subsection*{5. Design Science Research}

Für die methodische Einbettung der Arbeit. Priorität: \textbf{mittel}.

\begin{description}
  \item[Grundlagentext] \textcite{hevner2004dsr}: \emph{„Design Science in Information Systems Research"}, MIS Quarterly 28(1), 2004. Besonders die drei Design-Science-Zyklen (\emph{Relevance Cycle}, \emph{Design Cycle}, \emph{Rigor Cycle}) internalisieren; das ist das Gerüst für Kapitel 5 der Arbeit.

  \item[Praxisorientierter Einstieg] Peffers, K. et al.\ (2007): \emph{„A Design Science Research Methodology for Information Systems Research"}, JMIS 24(3). Gibt eine konkrete sechsstufige Prozessstruktur, die sich direkt auf diese Arbeit mappen lässt.
\end{description}

% -----------------------------------------------------------------------
\subsection*{6. Empfohlene GitHub-Repositories zum Erkunden}

\begin{description}
  \item[\texttt{tree-sitter/tree-sitter}] \url{https://github.com/tree-sitter/tree-sitter}: Hauptrepository; Grammar-Dateien für Python und TypeScript.
  \item[\texttt{modelcontextprotocol/python-sdk}] \url{https://github.com/modelcontextprotocol/python-sdk}: MCP Python-SDK mit Beispielen.
  \item[SWE-agent] \url{https://github.com/SWE-agent/SWE-agent}: Referenzimplementierung reaktiver Agent-Computer-Interfaces; für das Verständnis von B1.
  \item[Agentless] \url{https://github.com/OpenAutoCoder/Agentless}: Referenzimplementierung prozeduraler Kontextzusammenstellung; für das Verständnis des proaktiven Ansatzes.
\end{description}

% -----------------------------------------------------------------------
\subsection*{7. Abgeschlossene erste praktische Schritte}

\begin{enumerate}
  \item tree-sitter installieren und das WebshopTemplate parsen: Alle TypeScript-Funktionen und ihre Imports als JSON ausgeben.
  \item Diesen Output in eine Neo4j-AuraDB-Instanz schreiben: Knoten für Dateien und Funktionen, Kanten für IMPORTS.
  \item Eine Cypher-Abfrage schreiben: \emph{„Welche Funktionen werden von \texttt{CartContext.tsx} aufgerufen?"} Das ist der Proof-of-Concept für den Kern des Systems.
  \item Einen minimalen MCP-Server bauen, der diese Cypher-Abfrage als Tool bereitstellt, und ihn in Claude Code registrieren.
\end{enumerate}

Diese vier Schritte wurden im Rahmen der Einarbeitungsphase abgeschlossen und belegen, dass die technischen Kernkomponenten zusammenspielen. Sie bilden die Grundlage des in Abschnitt~4 beschriebenen Proof-of-Concept.


% ---------------------------------------------------------------------------
% Literaturverzeichnis
% ---------------------------------------------------------------------------
\newpage
\printbibliography[heading=bibintoc, title={Literaturverzeichnis}]

\end{document}
